\chapter{Les différents types de systèmes de fichiers}

Les systèmes de fichiers se classent principalement selon la nature du support sur lequel ils sont déployés et la manière dont ils gèrent la persistance et la sécurité des données.

\section{Systèmes de fichiers pour disques (Locaux)}
Ces systèmes sont conçus pour structurer le stockage sur des blocs physiques de données (disques durs, SSD, clés USB).

\begin{description}
    \item[Minix] Historiquement, le tout premier système de fichiers utilisé par Linux, hérité du système d'exploitation Minix. Il possédait de fortes limitations (noms de fichiers courts, taille limitée).
    \item[ext2 (Second Extended Filesystem)] Devenu pendant longtemps le standard de Linux. Il offre d'excellentes performances mais souffre d'une vulnérabilité importante en cas de coupure de courant soudaine (absence de journalisation).
    \item[msdos] Représente l'ancien système FAT16 standard utilisé à l'époque de MS-DOS et des premières versions de Windows. 
    \item[vfat] Extension du format FAT (FAT32) introduite avec Windows 95 pour lever la barrière des noms de fichiers courts. Il reste universellement compatible sur presque tous les appareils (clés USB, caméras).
    \item[ntfs (New Technology File System)] Le système de fichiers moderne par défaut de Windows (depuis Windows NT). Il intègre des fonctionnalités avancées comme la journalisation, la gestion des droits d'accès (permissions) et la compression.
\end{description}


\section{Systèmes de fichiers réseau (Partagés)}
Contrairement aux systèmes locaux, ils ne gèrent pas directement l'espace magnétique ou flash d'un disque dur local. Ils agissent comme une passerelle logicielle permettant d'accéder à distance aux fichiers stockés sur une autre machine à travers un protocole réseau.

\begin{description}
    \item[SMB / CIFS (Server Message Block)] Protocole de partage de fichiers développé par Microsoft, principalement utilisé dans les environnements Windows pour le partage de fichiers et d'imprimantes en réseau local.
    \item[NFS (Network File System)] Système de fichiers réseau traditionnel de l'univers Unix et Linux, initialement développé par Sun Microsystems. Il permet d'intégrer un répertoire distant de manière transparente dans l'arborescence locale.
\end{description}


\section{Systèmes de fichiers optiques (Médias amovibles)}
Ces formats sont standardisés pour garantir qu'un média optique gravé puisse être lu sur n'importe quel appareil ou système d'exploitation (Universalité).

\begin{description}
    \item[ISO 9660] Norme internationale définissant la structure de fichiers pour les CD-ROM. Elle garantit une compatibilité totale entre Windows, Mac et Unix.
    \item[UDF (Universal Disk Format)] Successeur de l'ISO 9660, optimisé pour les médias de plus grande capacité comme les DVD, Blu-ray et les supports réinscriptibles.
\end{description}


\section{Les systèmes de fichiers journalisés (Fiabilité)}
Un crash système ou une panne de courant pendant une écriture peut corrompre le système de fichiers.

La \textbf{journalisation} résout ce problème : avant de modifier un fichier, le système écrit l'action planifiée dans une zone réservée appelée \textbf{journal}.

\begin{description}
    \item[Le mécanisme] En cas de coupure brutale, le système n'a plus besoin de scanner l'intégralité du disque au redémarrage (opération très longue via \texttt{fsck} ou \texttt{chkdsk}). Il lui suffit de relire le journal pour finaliser les opérations interrompues ou annuler celles qui sont incomplètes, garantissant un retour rapide à un état cohérent.
\end{description}

Exemples de systèmes locaux intégrant nativement la journalisation :
\begin{itemize}[label=$\bullet$]
    \item \textbf{ext3 / ext4} : Évolutions journalisées d'ext2 sous Linux (\texttt{ext4} étant le standard moderne actuel).
    \item \textbf{ntfs} : Le système Windows utilise sa propre table journalisée pour sécuriser ses métadonnées.
    \item \textbf{ReiserFS / XFS} : Systèmes de fichiers alternatifs sous Linux, particulièrement performants pour la gestion des grands volumes de données ou de fichiers massifs (notamment \texttt{XFS}).
\end{itemize}